% ======================================================================
% SÉLECTION DE CONTENUS — Candidature MSI (Loupiac / Sage / Bonal / De Roland)
% ======================================================================
% Ce fichier NE CONTIENT AUCUNE MISE EN PAGE. Il choisit juste, parmi les
% macros de cv_banque_contenus.tex, lesquelles afficher et dans quel ordre.
%
% Ciblage : un seul CV pour candidater sur les 4 thématiques qui
% t'intéressent (pas un CV par thématique), en couvrant leur terrain commun
% plutôt qu'en sur-spécialisant sur une seule :
%   - Eric LOUPIAC   : Conception machines spéciales / Prototypage / Réorg. industrielle
%   - Louis SAGE     : Traitement d'images / informatique industrielle
%   - Mathieu BONNAL : Conception / Prototypage / Organisation industrielle / Gestion de flux
%   - Daniel DE ROLAND : Systèmes embarqués et objets connectés
%
% Rappel des blocs disponibles dans la banque :
%   Expérience : \expMicrotec \expPliageMetal \expGabarit
%   Associatif : \assoHandball \assoBenevole \assoGabarit
%   Projets    : \projBrasRobotique \projBuckConverter \projSatellite
%                \projBBEight \projQueueIA \projSoccerKit \projAppMobileAsso
%                \projNuitICAM \projExperiment \projGabarit
%   Profil     : \ProfilIngenierieSysteme \ProfilEmbarqueLogiciel
%                \ProfilGeneraliste (ou un texte libre)
%
% Convention : chaque macro de sélection suppose qu'une coordonnée de départ
% (exp_pos / asso_pos / proj_pos) a déjà été positionnée par le CV principal
% juste avant l'appel (sous le titre de la rubrique).
% ======================================================================

% --- PROFIL ---
% Version généraliste : pas de candidature MSI mono-thématique, donc pas
% de profil sur-ciblé aéro/systèmes (garde ça pour MBDA).
\newcommand{\Profil}{\ProfilGeneraliste}

% --- COMPÉTENCES ---
\newcommand{\Competences}{%
    \item Rigueur Scientifique
    \item Esprit d'équipe
    \item Adaptabilité
    \item Polyvalence
    \item Autonomie
    \item Gestion de projet
}

% --- OUTILS & LANGAGES ---
% Conservé large : SolidWorks/CAO et Arduino/ESP32/Linux couvrent
% conception-prototypage + embarqué/IoT ; Python couvre l'informatique
% industrielle et le traitement de données.
\newcommand{\OutilsLangages}{%
    \item Python (Algorithmie)
    \item C / C++ (Embarqué)
    \item Matlab / Simulink
    \item Arduino / ESP32
    \item SolidWorks / CAO
    \item Git / GitHub
    \item ROS2 / RoboDK
}

% --- EXPÉRIENCE ---
% Les deux expériences se complètent bien ici : Microtec pour le bureau
% d'études, Pliage Métal pour l'atelier/la production (parle directement à
% Loupiac et Bonnal sur l'organisation industrielle).
\newcommand{\SelectionExperience}{%
    \expMicrotec{exp_pos}
    \chainNext{exp_pos}{exp_last}{\gapTextToSub}
    \expPliageMetal{exp_pos}
}

% --- ENGAGEMENT ASSOCIATIF ---
\newcommand{\SelectionAssociatif}{%
    \assoHandball{asso_pos}
    \chainNext{asso_pos}{asso_last}{\gapTextToSub}
    \assoBenevole{asso_pos}
}

% --- PROJETS RÉALISÉS ---
% 4 projets choisis pour couvrir ENSEMBLE les 4 thématiques visées, sans
% qu'aucun ne soit "hors sujet" pour l'un des 4 encadrants :
%   - Bras Robotique   : conception/prototypage (Loupiac, Bonnal) + embarqué (De Roland)
%   - Queue IA         : traitement d'images/vision (Sage) + objets connectés (De Roland)
%   - BB8              : programmation embarquée / asservissement (De Roland)
%   - Nuit ICAM        : organisation industrielle / gestion de flux à grande échelle (Loupiac, Bonnal)
\newcommand{\SelectionProjets}{%
    \projBrasRobotique{proj_pos}
    \chainNext{proj_pos}{proj_last}{0.05cm}
    \projQueueIA{proj_pos}
    \chainNext{proj_pos}{proj_last}{0.05cm}
    \projBBEight{proj_pos}
    \chainNext{proj_pos}{proj_last}{0.05cm}
    \projNuitICAM{proj_pos}
}
